% Canonical TeX edition. Initial text inherited from 3b067cc2eba4ac2f2f91248448727639d8ac7c5a
% Source: research/case_study/phase238_L083_application.md
% Subsequent mathematical edits require a new audit.
% Options for packages loaded elsewhere
\PassOptionsToPackage{unicode}{hyperref}
\PassOptionsToPackage{hyphens}{url}
\documentclass[
  11pt,
  a4paper,
]{article}
\usepackage{xcolor}
\usepackage[margin=1in]{geometry}
\usepackage{amsmath,amssymb}
\setcounter{secnumdepth}{-\maxdimen} % remove section numbering
\usepackage{iftex}
\ifPDFTeX
  \usepackage[T1]{fontenc}
  \usepackage[utf8]{inputenc}
  \usepackage{textcomp} % provide euro and other symbols
\else % if luatex or xetex
  \usepackage{unicode-math} % this also loads fontspec
  \defaultfontfeatures{Scale=MatchLowercase}
  \defaultfontfeatures[\rmfamily]{Ligatures=TeX,Scale=1}
\fi
\usepackage{lmodern}
\ifPDFTeX\else
  % xetex/luatex font selection
  \setmainfont[]{STIXGeneral}
  \setmonofont[]{Menlo}
\fi
% Use upquote if available, for straight quotes in verbatim environments
\IfFileExists{upquote.sty}{\usepackage{upquote}}{}
\IfFileExists{microtype.sty}{% use microtype if available
  \usepackage[]{microtype}
  \UseMicrotypeSet[protrusion]{basicmath} % disable protrusion for tt fonts
}{}
\makeatletter
\@ifundefined{KOMAClassName}{% if non-KOMA class
  \IfFileExists{parskip.sty}{%
    \usepackage{parskip}
  }{% else
    \setlength{\parindent}{0pt}
    \setlength{\parskip}{6pt plus 2pt minus 1pt}}
}{% if KOMA class
  \KOMAoptions{parskip=half}}
\makeatother
\usepackage{longtable,booktabs,array}
\newcounter{none} % for unnumbered tables
\usepackage{calc} % for calculating minipage widths
% Correct order of tables after \paragraph or \subparagraph
\usepackage{etoolbox}
\makeatletter
\patchcmd\longtable{\par}{\if@noskipsec\mbox{}\fi\par}{}{}
\makeatother
% Allow footnotes in longtable head/foot
\IfFileExists{footnotehyper.sty}{\usepackage{footnotehyper}}{\usepackage{footnote}}
\makesavenoteenv{longtable}
\setlength{\emergencystretch}{3em} % prevent overfull lines
\providecommand{\tightlist}{%
  \setlength{\itemsep}{0pt}\setlength{\parskip}{0pt}}
\usepackage{newunicodechar}
\newfontfamily\numeralfallback{STIX Two Text}
\newunicodechar{∎}{\ensuremath{\blacksquare}}
\usepackage{fvextra}
\DefineVerbatimEnvironment{Highlighting}{Verbatim}{breaklines,breakanywhere,breakanywheresymbolpre={},breakanywheresymbolpost={},breaksymbolleft={},breaksymbolright={},commandchars=\\\{\}}
\newunicodechar{⁰}{{\numeralfallback\char"2070\relax}}
\newunicodechar{₀}{{\numeralfallback\char"2080\relax}}
\newunicodechar{₁}{{\numeralfallback\char"2081\relax}}
\newunicodechar{₂}{{\numeralfallback\char"2082\relax}}
\newunicodechar{₃}{{\numeralfallback\char"2083\relax}}
\newunicodechar{₄}{{\numeralfallback\char"2084\relax}}

\usepackage{seqsplit}
\renewcommand{\texttt}[1]{{\ttfamily\seqsplit{#1}}}
\setlength{\emergencystretch}{6em}

\newunicodechar{⁻}{{\numeralfallback\char"207B\relax}}
\newunicodechar{⁸}{{\numeralfallback\char"2078\relax}}
\newunicodechar{ⁱ}{{\numeralfallback\char"2071\relax}}
\newunicodechar{⁴}{{\numeralfallback\char"2074\relax}}
\newunicodechar{⁷}{{\numeralfallback\char"2077\relax}}
\newunicodechar{⁶}{{\numeralfallback\char"2076\relax}}
\newunicodechar{⁵}{{\numeralfallback\char"2075\relax}}
\newunicodechar{₈}{{\numeralfallback\char"2088\relax}}
\newunicodechar{ᵢ}{\textsubscript{i}}
\newunicodechar{ⱼ}{\textsubscript{j}}
\newunicodechar{ₜ}{\textsubscript{t}}
\usepackage{bookmark}
\IfFileExists{xurl.sty}{\usepackage{xurl}}{} % add URL line breaks if available
\urlstyle{same}
% fallback for those not using the hyperref driver hyperxmp:
\makeatletter
\@ifundefined{xmpquote}{\newcommand{\xmpquote}[1]{#1}}{}
\makeatother
\hypersetup{
  hidelinks,
  pdfcreator={LaTeX via pandoc}}

\author{}
\date{}

\begin{document}

\section{The two projective special linear
four-polytopes}\label{the-two-projective-special-linear-four-polytopes}

\subsection{Abstract}\label{abstract}

We derive the classification of regular polytopes of rank at least four
with full automorphism group \(\operatorname{PSL}_2(q)\) from the common
projective-line reduction. The only examples are Grünbaum's 11-cell and
Coxeter's 57-cell. The proof includes uniqueness of the marked groups,
their exact universal presentations, self-duality, face counts and the
covering universality of their local maps. The classification is due to
Leemans and Schulte; the purpose of this paper is its complete
application of the shared centralizer and finite-amalgam theory.

\subsection{1. Statement and the common
theorem}\label{statement-and-the-common-theorem}

Let q be a prime power. We use the string C-group convention of the
companion paper {[}E{]}: the distinguished generators are nonidentity
involutions, they commute at nonconsecutive positions, and every pair of
marked subgroups has the prescribed intersection. We retain the standard
correspondence between these marked groups and regular abstract
polytopes {[}MS{]}. Group isomorphisms preserving the ordered marking
give polytope isomorphisms; reversal gives duality. The automorphism
groups below are the full flag-regular groups.

\textbf{Theorem 1 (Leemans--Schulte).} A regular abstract polytope of
rank at least four whose full automorphism group is isomorphic to
\(\operatorname{PSL}_2(q)\) has rank four. It is precisely one of the
following two polytopes:

{\def\LTcaptype{none} % do not increment counter
\begin{longtable}[]{@{}rllll@{}}
\toprule\noalign{}
q & Polytope & Type & Facets & Vertex figures \\
\midrule\noalign{}
\endhead
\bottomrule\noalign{}
\endlastfoot
11 & 11-cell & \{3, 5, 3\} & hemi-icosahedra & hemi-dodecahedra \\
19 & 57-cell & \{5, 3, 5\} & hemi-dodecahedra & hemi-icosahedra \\
\end{longtable}
}

Each row gives one marked-group and polytope isomorphism class. Both
polytopes are self-dual and universal among regular polytopes with the
indicated facets and vertex figures. Their face vectors, from vertices
through facets, are (11, 55, 55, 11) and (57, 171, 171, 57),
respectively. The 11-cell has complete edge graph \(K_{11}\). Both full
groups equal their even-generator subgroups.

The common theorem {[}E, Theorem 1{]} applies to the entire prime-power
range. It gives odd characteristic, rank four, and facet and
vertex-figure groups \(S_4\) or \(A_5\). It then supplies all six
ordered rank-three marking classes and the ten possible universal
presentations, up to reversal. This uses the retained Dickson subgroup
types, the proved centralizer descent and odd-order normalizer argument,
and the exact finite computation of {[}E, §5{]}. In particular, it
already supplies the exclusion of small characteristic-two groups; no
unproved rank-three existence assertion is used here.

\subsection{2. Exhaustion of the possible
quotients}\label{exhaustion-of-the-possible-quotients}

Suppose that \(G=\operatorname{PSL}_2(q)\) has the required marking. The
two end parabolics are proper by IP. We apply the ten-row table of {[}E,
Lemma 11{]} with its literal markings. Every possible G is a marked
quotient of one of those universal groups. It is necessary to examine
quotients as well as the groups themselves.

The four trivial rows give no marking by nonidentity involutions. In the
rows of order 60 and 24, an end parabolic equals the entire universal
group. Thus the missing endpoint generator is a word in the other three.
This same dependence survives in every quotient, contradicting IP. These
two rows are excluded even before using the ambient projective group.

In the row of order 96, the order of G must divide 96. For odd
\(q\ge 5\), the smallest value of \(q(q^2-1)/2\) is 60, which does not
divide 96; at \(q\ge 7\) the order is already at least 168. The
remaining odd value \(q=3\) gives \(A_4\), which cannot contain a proper
facet \(S_4\) or \(A_5\). The even values were excluded by the common
rank theorem. Hence the order 96 row contributes nothing.

The row of order 120 is \(S_5\). Its normal subgroups are 1, \(A_5\),
\(S_5\), so its only proper nontrivial quotient is \(C_2\). Neither
\(C_2\) nor \(S_5\) is one of the possible odd-characteristic
\(\operatorname{PSL}_2(q)\): for \(q\ge 5\) the latter is nonabelian
simple, and the remaining small group is \(A_4\). Thus this row also
gives nothing.

The remaining universal groups are \(\operatorname{PSL}_2(11)\), of
order 660, and \(\operatorname{PSL}_2(19)\), of order 3420. Both are
simple. Any nontrivial marked quotient is therefore the whole group. For
odd q the formula \(q(q^2-1)/2\) is strictly increasing, so the two
possibilities force \(q=11\) and \(q=19\). This proves exhaustion for
every prime power.

Conversely, the exact presentations in those two rows have the displayed
groups and faithful distinguished markings. The common finite proof
verifies all subset-intersection identities, as well as generation and
the generator orders. Their consecutive labels are exactly 3, 5, 3 and
5, 3, 5: the presentation bounds the orders by these primes, and the
distinct involutions prevent order 1. The correspondence {[}MS{]}
therefore gives regular rank-four polytopes with exactly those full
automorphism groups and types.

Uniqueness follows at the level required by the correspondence. Each
local marking lies in one of the six complete classes in {[}E, Lemma
10{]}, and each compatible pair has its exact universal presentation.
The only surviving universal groups are simple, so every admissible
quotient map is an isomorphism preserving all four marked generators.
The common reversal automorphisms identify the two orientations of each
local pair. Consequently there is one polytope in each row, not merely
one abstract group. This proves the classification and uniqueness
assertions of Theorem 1.

\subsection{3. Local maps, geometry and
universality}\label{local-maps-geometry-and-universality}

For clarity, the two exact presentations consist of the ordinary string
Coxeter relations with type \{3, 5, 3\} or \{5, 3, 5\}, together with

\[
 (r_0r_1r_2)^5=(r_1r_2r_3)^5=1.
 \tag{2}
\]

There are no extra relators hidden in this description: these are
precisely the two surviving presentations in the complete common table.
Their local presentations are \([3,5]_5\) and \([5,3]_5\) in the Petrie
convention of {[}E, equation (5){]}. The corresponding standard regular
maps are the antipodal quotients of the icosahedron and dodecahedron.
This classical identification supplies the names hemi-icosahedron and
hemi-dodecahedron and their projective-plane interpretation. It is
retained from the geometric conventions in {[}LS, §5{]} and {[}MS{]}.
Thus the polytopes are locally projective, with the facets and vertex
figures stated in the theorem.

Each facet and vertex figure has group \(A_5\) of order 60. Hence the
numbers of facets and vertices are 660/60=11 and 3420/60=57. If the
common endpoint label is m, the stabilizer of an edge is
\(C_2\times D_{2m}\), of order 4m: the isolated endpoint generator
commutes with the opposite adjacent pair, and IP makes their
intersection trivial. The two-face stabilizer has the same order. Thus
there are 660/12=55 edges and triangular two-faces in the first case,
and 3420/20=171 edges and pentagonal two-faces in the second. These are
the face vectors in the theorem. The actual edge-coset calculation of
{[}E, Lemma 11{]} identifies 55 distinct unordered endpoint pairs in the
first case; there are exactly 55 pairs on 11 vertices. The edge graph is
therefore \(K_{11}\), not an inference from the face numbers alone.

The reversal homomorphisms in {[}E, Lemma 11{]} are automorphisms of the
actual marked groups. By {[}MS{]} they give the claimed self-dualities.
The subgroup of even-length words in the distinguished generators has
index at most two in any group generated by involutions. Our two groups
are nonabelian simple, so this subgroup is the whole group.

Finally let Q be any regular polytope with one of the stated facet and
vertex-figure pairs. Its distinguished generators satisfy the
corresponding universal presentation above. They therefore receive a
marked epimorphism from the already verified string C-group in the
table. The retained marked-epimorphism/covering theorem {[}MS{]} gives a
cover from our polytope to Q. This proves the asserted universality. The
argument does not assume that a general abstract group amalgam satisfies
IP; that property was proved for these exact universal groups. It does
not require Q to be assumed finite, since the epimorphism itself forces
its group to be finite. The remaining assertions of Theorem 1 follow. ∎

\subsection{4. Attribution and verification
boundary}\label{attribution-and-verification-boundary}

The source contract is {[}LS, Theorems 1, 3, 4{]}, its introductory rank
restriction, its complete ten-row Table 1 and the specified geometric
properties in §5. The 11-cell and 57-cell are the prior constructions of
Grünbaum and Coxeter. The classification belongs to Leemans and Schulte.
The present application uses the common replacement proof of {[}E{]} and
supplies the complete projective-special-linear quotient and geometry
arguments above.

The exact finite computation and its program belong to the shared paper,
not to an additional computation repeated by this application. Its five
retained Lean declarations certify the abstract odd-order normalizer
obstruction. They do not certify this whole application, the matrix
descent, the finite universal orders, or the polytope correspondence.
The required GAP computation remains an explicit computational
dependency.

The historical subfield-intersection lemmas and their multiplicity
formulas are unnecessary to this proof. The source's one-Petrie-relation
groups and its later discussion of two-flag-orbit constructions are
outside the present contract. No new open-question solution or
historical-priority claim is made.

\subsection{References}\label{references}

{[}E{]} \emph{Centralizers, facets and finite amalgams on the projective
line}, companion umbrella paper, especially Theorem 1 and Lemmas 10--11.

{[}LS{]} D. Leemans and E. Schulte, \emph{Groups of type \(L_2(q)\)
acting on polytopes}, retained arXiv:math/0606660v1 author edition, 14
pages; journal publication \emph{Advances in Geometry} 7 (2007),
529--539. The exact primary edition is bound in the accompanying source
contract.

{[}MS{]} P. McMullen and E. Schulte, \emph{Abstract Regular Polytopes},
Cambridge University Press, 2002, especially the string C-group and
covering correspondence.

\end{document}
